\documentclass[11pt]{article}
% Shared preamble for all daily exercises
% Update this file to change settings (padding, parindent, etc.) for all exercises

\usepackage[margin=0.75in]{geometry}
\usepackage{amsmath}
\usepackage{amssymb}

% Custom settings
\setlength{\parindent}{0pt}
\setlength{\parskip}{0.2cm}

\title{Daily Exercise}
\author{Dhruman Gupta}
\date{19th April}

\begin{document}

% Common header for daily exercises
% This file can be included in other LaTeX documents

\maketitle

\noindent
\textbf{Course:} CS-3330, Theory of Computation

\vspace{0.2cm}

\noindent
\textbf{Question:}\noindent 
Show that there exists an oracle $A$ such that $P^A = NP^A$. Show that there exists an oracle $B$ such that $P^B \neq NP^B$.

\textbf{Answer:}

First, we show there exists an oracle $A$ such that $P^A = NP^A$.

Let $A = \text{EXPCOM}$, where $\text{EXPCOM}$ is some EXP-complete language.

We claim that
\[
P^A = EXP.
\]

To see that $EXP \subseteq P^A$, let $L \in EXP$. Since $A$ is EXP-complete, there exists a polynomial time many-one reduction $f$ such that
\[
x \in L \iff f(x) \in A.
\]
So a deterministic polynomial time machine with oracle access to $A$ can decide $L$ by computing $f(x)$ and making one oracle query. Hence $L \in P^A$.

Now, to see that $P^A \subseteq EXP$, note that a polynomial time oracle machine can only make polynomially many queries, each of polynomial length. Since $A \in EXP$, each oracle query can be answered in exponential time. Therefore the whole oracle computation can be simulated in exponential time. Hence $P^A \subseteq EXP$.

So $P^A = EXP$.

Now consider $NP^A$. We claim that
\[
NP^A \subseteq EXP.
\]
Indeed, an $NP^A$ machine runs in polynomial time, so it has at most exponentially many computation branches. On each branch it makes polynomially many oracle queries of polynomial length. Since each query to $A$ can be answered in exponential time, the entire nondeterministic computation can be simulated deterministically in exponential time. Hence $NP^A \subseteq EXP$.

But we already showed that $P^A = EXP$, so
\[
NP^A \subseteq EXP = P^A.
\]
Also, clearly $P^A \subseteq NP^A$. Therefore
\[
P^A = NP^A.
\]

Now we show there exists an oracle $B$ such that $P^B \neq NP^B$.

Define
\[
L_B = \{1^n \mid \exists y \in \{0,1\}^n \text{ such that } y \in B\}.
\]
First, note that $L_B \in NP^B$. On input $1^n$, an NP machine guesses a string $y \in \{0,1\}^n$ and queries whether $y \in B$. If yes, it accepts. Otherwise it rejects. So $L_B \in NP^B$.

Now we construct $B$ so that no deterministic polynomial time oracle machine decides $L_B$.

Enumerate all deterministic polynomial time oracle Turing machines as
\[
M_1, M_2, M_3, \dots
\]
and let $p_i$ be a polynomial time bound for $M_i$.

We build $B$ in stages. At stage $i$, choose a number $n$ such that:
\begin{itemize}
    \item $n$ is larger than every length used in earlier stages,
    \item $p_i(n) < 2^n$.
\end{itemize}

We will decide membership in $B$ only for strings of length $n$ at this stage. All answers for shorter lengths were already fixed in earlier stages. For strings of length $n$, temporarily answer every query by saying ``no''.

Now simulate $M_i$ on input $1^n$ with these oracle answers. Since $M_i$ runs in time at most $p_i(n)$, it asks at most $p_i(n)$ queries in total. In particular, it can ask at most $p_i(n)$ strings of length $n$. But there are $2^n$ such strings, and $p_i(n) < 2^n$, so there is some string $z \in \{0,1\}^n$ that $M_i$ never queries.

Now diagonalize against $M_i$:
\begin{itemize}
    \item If $M_i$ accepts $1^n$, then put no string of length $n$ into $B$.
    \item If $M_i$ rejects $1^n$, then put $z$ into $B$.
\end{itemize}

This guarantees that $M_i$ is wrong on input $1^n$.

If $M_i$ accepts, then no string of length $n$ is in $B$, so $1^n \notin L_B$, but $M_i$ accepts.

If $M_i$ rejects, then $z \in B$, so there exists a string of length $n$ in $B$. Hence $1^n \in L_B$, but $M_i$ rejects.

Also, in the second case, adding $z$ to $B$ does not change the computation of $M_i$, because $M_i$ never queried $z$.

So stage $i$ makes sure that $M_i$ does not decide $L_B$. Since this is true for every deterministic polynomial time oracle machine, no machine in $P^B$ decides $L_B$. Thus
\[
L_B \notin P^B.
\]
But we already showed that $L_B \in NP^B$. Hence
\[
P^B \neq NP^B.
\]

Therefore there exists an oracle $A$ such that $P^A = NP^A$, and there exists an oracle $B$ such that $P^B \neq NP^B$.

\end{document}
